% Source: https://artofproblemsolving.com/community/c3h2031599_midpoint_type_formula_if_you_want_to_find_other_points_on_a_line_segment
% Stable identity: aops-topic-2031599-q01
(b) Going from $A$ to $B$ changes $x$ by $x_2-x_1$ and $y$ by $y_2-y_1$. Moving along the segment, the changes in $x$ and $y$ grow in proportion to the distance travelled (similar right triangles with legs parallel to the axes), so the point $t$ of the way along is
\[
\bigl(x_1+t(x_2-x_1),\ y_1+t(y_2-y_1)\bigr)=\bigl((1-t)x_1+tx_2,\ (1-t)y_1+ty_2\bigr).
\]
Its distance from $A$ is $t\sqrt{(x_2-x_1)^2+(y_2-y_1)^2}=t\cdot AB$, as required. For $t=\dfrac{1}{2}$ this is the midpoint formula.

(a) Here the run is $9-2=7$ and the rise is $4-3=1$.
\begin{asy}
usepackage("amsmath");
size(7cm);
defaultpen(fontsize(10pt));
pair A=(2,3), B=(9,4), P=(24/5,17/5);
draw((0,0)--(10,0),gray,Arrow(5)); draw((0,0)--(0,5.2),gray,Arrow(5));
draw(A--B,linewidth(0.9));
draw(A--(9,3),dashed+gray); draw((9,3)--B,dashed+gray);
dot(A); dot(B); dot(P,red);
label("$A(2,3)$",A,W); label("$B(9,4)$",B,E);
label("$P\left(\dfrac{24}{5},\dfrac{17}{5}\right)$",P,N,red);
label("$7$",(5.5,3),S); label("$1$",(9,3.5),E);
label("\mbox{}",(min(currentpicture,user=true)+max(currentpicture,user=true))/2,black); // keeps the text after the figure black when LaTeX typesets the labels inline
\end{asy}

With $t=\dfrac{2}{5}$,
\[
P=\left(2+\dfrac{2}{5}\cdot7,\ 3+\dfrac{2}{5}\cdot1\right)=\boxed{\left(\dfrac{24}{5},\dfrac{17}{5}\right)}.
\]
